\section*{Chapitre \ref{ch:ml:from-prediction-to-portfolio} --- De la prédiction au portefeuille}

\begin{solution}{exo:ml:from-prediction-to-portfolio:1}
Maximiser $\mu w - \frac\gamma2\sigma^2w^2 - \kappa(w - w_{\mathrm{prev}})^2$~: la condition du premier ordre $\mu - \gamma\sigma^2w -
2\kappa(w - w_{\mathrm{prev}}) = 0$ donne la règle. Avec $\kappa = 0$, on obtient le poids moyenne-variance $\mu/(\gamma\sigma^2)$~; lorsque
$\kappa\to\infty$, le poids reste égal à $w_{\mathrm{prev}}$.
\end{solution}

\begin{solution}{exo:ml:from-prediction-to-portfolio:2}
La règle est $w = a\,\mu/(\gamma\sigma^2) + (1 - a)w_{\mathrm{prev}}$ avec $a = \gamma\sigma^2/(\gamma\sigma^2 + 2\kappa) =
0.018/(0.018 + 0.001) = 0.947$~: 95~\% de l'écart est refermé à chaque période~; l'approximation du coût ne ralentit le portefeuille que légèrement.
\end{solution}

\begin{solution}{exo:ml:from-prediction-to-portfolio:3}
Multiplier chaque rendement par une constante laisse le ratio de Sharpe inchangé, si bien que toute échelle des coefficients de même
sens est tout aussi bonne, aux effets de coûts près~; la montée de gradient dérive en échelle, des positions plus importantes génèrent davantage de transactions, et le
turnover de l'\omterm{def:ml:from-prediction-to-portfolio:dfl}{apprentissage orienté décision} (2.75) dépasse celui des moindres carrés alignés (1.59) pour le même ratio de Sharpe.
\end{solution}

\begin{solution}{exo:ml:from-prediction-to-portfolio:4}
L'erreur quadratique est dominée par les titres non négociables, où le signal de la première caractéristique est fort~; l'ajuster réduit
davantage l'erreur à cet endroit qu'il ne l'augmente parmi les titres négociables, où ce coefficient n'est qu'une erreur pure. Les moindres carrés sur les
titres négociables ne s'ajustent que là où le portefeuille l'évalue.
\end{solution}

\begin{solution}{exo:ml:from-prediction-to-portfolio:5}
Lorsque la sensibilité du portefeuille à la prévision varie d'un titre négociable ou d'une période à l'autre d'une façon qu'un poids ne peut exprimer~:
des limites de position contraignantes pour certains titres, des coûts qui dépendent de la taille des transactions, un modèle de risque factoriel qui compense les positions, une
prévision qui alimente une règle non linéaire. Alors, les erreurs qui comptent ne sont pas simplement celles portant sur les titres négociés.
\end{solution}

\begin{solution}{exo:ml:from-prediction-to-portfolio:6}
La comparaison se fait dans le backtest sur lequel le modèle de bout en bout a été ajusté ou sélectionné~: une fonction de perte définie sur le ratio de Sharpe du portefeuille
est une forte incitation à ajuster le bruit. Comparer sur des données non utilisées pour l'ajustement, face à une prévision entraînée sur les titres négociés, avec la dispersion
sur des échantillons indépendants~; le cas à petit échantillon du chapitre montre que l'ajustement de bout en bout perd hors échantillon.
\end{solution}

\begin{solution}{exo:ml:from-prediction-to-portfolio:7}
Partant des moindres carrés sur les titres négociables, avec le nombre d'\omterm{def:ml:neural-networks-for-noisy-tabular-data:backprop}{époques} choisi sur les 60 dernières périodes d'entraînement, le
bloc de validation ne retient aucun entraînement de bout en bout dans trois des cinq échantillons~; le ratio de Sharpe net moyen en test est de 1.04,
le même que celui des moindres carrés alignés~. L'\omterm{def:ml:trees-and-boosting:early}{arrêt précoce} comble l'écart surtout en s'arrêtant dès le départ.
\end{solution}

\begin{solution}{exo:ml:from-prediction-to-portfolio:8}
$w^* = \Sigma^{-1}\mu/\gamma$, donc $\partial w^*/\partial\mu = \Sigma^{-1}/\gamma$. Pour une perte $L(w^*)$, $\partial L/\partial\mu =
(\partial w^*/\partial\mu)^\top\nabla_wL = \Sigma^{-1}\nabla_wL/\gamma$, et la \omterm{def:ml:neural-networks-for-noisy-tabular-data:backprop}{rétropropagation} la reporte jusqu'aux
paramètres de la prévision. L'erreur de la prévision compte dans les directions qu'amplifie l'optimiseur~: celles de faible risque.
\end{solution}

\begin{solution}{pb:ml:from-prediction-to-portfolio:1}
\textbf{Partie I.}
\begin{enumerate}
\item Titres non négociables où la première caractéristique prédit fortement, titres négociables où la seconde prédit faiblement~; une
seule prévision linéaire ne peut avoir raison dans les deux cas.
\item Poids moyenne-variance avec un proxy quadratique de coût sur les titres négociables~; les rendements nets acquittent 5 points de base sur chaque variation
de poids.
\item 1.51 sur la première caractéristique et 0.43 sur la seconde (unités de 0.1~\%)~: surtout le signal des titres non négociables, qui domine
l'erreur quadratique.
\item Ensemble des titres~: 9.007 au total et 8.987 sur les titres négociables~; ajustement sur les titres négociables~: 9.031 et 8.962 (unités de $10^{-4}$).
\end{enumerate}
\textbf{Partie II.}
\begin{enumerate}[resume]
\item Net, brut, turnover~: 0.19, 0.60, 2.95~; 1.52, 1.92, 1.59~; 1.51, 1.92, 2.75~; 1.49, 1.92, 2.90~; vérité terrain 1.57, 1.97, 1.50.
\item 1.33, avec un écart-type de 0.18 sur cinq paires.
\item Les deux apprennent la direction de la seconde caractéristique~; ici, ce que la décision apporte se limite à savoir quels titres comptent, et un
ajustement restreint dit la même chose.
\item Elle n'a pas de lissage~: les poids suivent les caractéristiques à chaque période.
\end{enumerate}
\textbf{Partie III.}
\begin{enumerate}[resume]
\item Moindres carrés alignés 1.04, \omterm{def:ml:from-prediction-to-portfolio:dfl}{apprentissage orienté décision} 0.89, paramétrique 0.87, \omterm{def:ml:from-prediction-to-portfolio:pto}{prédire puis optimiser} 0.09, vérité terrain 1.51.
\item Les coefficients sur les caractéristiques inutiles, et un turnover de 3.49 contre 1.88.
\item Peu de paramètres, un démarrage par moindres carrés, un \omterm{def:ml:trees-and-boosting:early}{arrêt précoce} sur un bloc de validation, des pénalités sur la direction, et non seulement
sur l'échelle.
\item Sur des périodes non utilisées pour l'ajustement et des échantillons indépendants, selon la performance nette et sa dispersion, avec les moindres carrés alignés comme
indice de référence.
\end{enumerate}
\textbf{Partie IV.}
\begin{enumerate}[resume]
\item \emph{La perte qui rapporte.} \omterm{def:ml:from-prediction-to-portfolio:pto}{Prédire puis optimiser} rapporte un ratio de Sharpe net de 0.19, l'\omterm{def:ml:from-prediction-to-portfolio:dfl}{apprentissage orienté décision} 1.51 et la
\omterm{def:ml:reinforcement-learning-foundations:mdp}{politique} paramétrique 1.49~; l'écart de 1.33 a un écart-type de 0.18 sur cinq paires entraînement-test, et les moindres carrés sur les
titres négociés seuls égalent les méthodes de bout en bout (1.52).
\item Lorsque le modèle de prévision est bien spécifié, ou lorsqu'il est ajusté sur les titres et les périodes qu'utilise la décision.
\item Décisions dont la sensibilité à la prévision varie d'une façon qu'une restriction d'échantillon ou un poids ne peut saisir, et assez de données
pour les ajuster.
\item Le coefficient de transfert mesure ce que les contraintes perdent entre la prévision et le portefeuille~; la perte de ce chapitre survient
en amont, lorsque la prévision apprend la mauvaise chose.
\item Sa performance nette de portefeuille selon la construction du desk et ses coûts, sur des données non utilisées pour l'ajustement, face au modèle en place, avec
le turnover et la fourchette.
\item Toujours dans l'évaluation du portefeuille~; à l'entraînement, par une règle tenant compte des coûts ou une perte pondérée par les coûts.
\item La même chose, avec le \omterm{def:ml:neural-networks-for-noisy-tabular-data:mlp}{réseau de neurones} à la place de la prévision linéaire, et davantage de régularisation.
\item Faire en sorte que le portefeuille construit à partir d'elle soit bon, net des coûts, hors échantillon.
\end{enumerate}
\end{solution}

\begin{solution}{iq:ml:from-prediction-to-portfolio:1}
La précision se mesure partout~; le profit ne dépend de la prévision que là où l'optimiseur s'appuie sur elle, et dans la mesure de cet appui, après
coûts et contraintes. Un modèle mal spécifié transfère des erreurs d'un titre à l'autre, et la perte décide lesquelles.
\iqlookfor{mauvaise spécification et décalage entre la perte et la décision.}
\end{solution}

\begin{solution}{iq:ml:from-prediction-to-portfolio:2}
Entraîner la prévision selon la perte de la décision prise à partir d'elle~; dériver la solution de l'optimiseur sous forme fermée, par
déroulage d'un solveur itératif, ou par différentiation implicite de ses conditions d'optimalité.
\iqlookfor{la perte et au moins une méthode de différentiation.}
\end{solution}

\begin{solution}{iq:ml:from-prediction-to-portfolio:3}
Les poids sont une fonction linéaire des caractéristiques autour d'un indice de référence, avec des coefficients choisis pour maximiser l'utilité moyenne
dans l'échantillon (in-sample)~. Peu de paramètres, pas de matrice de covariance~; les coûts et les contraintes peuvent entrer dans l'utilité, mais la maximisation
de l'utilité dans l'échantillon conduit au \omterm{def:ml:why-financial-machine-learning-is-different:generalisation}{surapprentissage}, et la forme de la \omterm{def:ml:reinforcement-learning-foundations:mdp}{politique} est restrictive.
\iqlookfor{la forme, le critère d'ajustement et le risque de \omterm{def:ml:why-financial-machine-learning-is-different:generalisation}{surapprentissage}.}
\end{solution}

\begin{solution}{iq:ml:from-prediction-to-portfolio:4}
Évaluer le modèle sur les rendements nets du portefeuille~; l'ajuster sur les titres négociables, pondérés par la liquidité ou la position attendue~; ou entraîner
par une règle de portefeuille tenant compte des coûts~; et pénaliser le turnover dans la prévision (lissage) là où les coûts sont élevés.
\iqlookfor{évaluation nette et au moins une façon d'intégrer les coûts à l'entraînement.}
\end{solution}

\begin{solution}{iq:ml:from-prediction-to-portfolio:5}
Si le modèle a été évalué sur des données utilisées pour l'ajustement ou la sélection, son turnover et ses coûts, la dispersion selon les graines et les échantillons,
l'indice de référence (une prévision ajustée sur les titres négociés) et le nombre de paramètres au regard des données.
\iqlookfor{évaluation hors échantillon, un indice de référence équitable et une dispersion.}
\end{solution}

\begin{solution}{iq:ml:from-prediction-to-portfolio:6}
Pour $\min\frac12w^\top Qw - \mu^\top w$ sous la contrainte $Aw = b$, le système KKT $Qw - \mu + A^\top\lambda = 0$, $Aw = b$ est linéaire~;
la différentiation donne $\begin{pmatrix}Q & A^\top\\ A & 0\end{pmatrix}\begin{pmatrix}dw\\ d\lambda\end{pmatrix} =
\begin{pmatrix}d\mu\\ 0\end{pmatrix}$, donc $dw/d\mu$ est le bloc supérieur gauche de l'inverse~; les contraintes d'inégalité ajoutent les
contraintes actives à $A$ (OptNet).
\iqlookfor{le système KKT linéarisé et les contraintes actives.}
\end{solution}
